% Préambule commun aux manuels de Tutti (français et anglais).
% Compilation : pdflatex manuel-fr.tex (deux fois), idem manuel-en.tex, depuis docs/src/.
\documentclass[11pt,a4paper]{article}
\usepackage[utf8]{inputenc}
\usepackage[T1]{fontenc}
\usepackage{lmodern}
\usepackage{microtype}
\usepackage[a4paper,margin=2cm,headheight=24pt]{geometry}
\usepackage{graphicx}
\usepackage[table]{xcolor}
\usepackage[most]{tcolorbox}
\usepackage{enumitem}
\usepackage{booktabs}
\usepackage{array}
\usepackage{tabularx}
\usepackage[nobottomtitles*]{titlesec}
\usepackage{needspace}
\usepackage{fancyhdr}
\usepackage{fontawesome5}
\usepackage{etoolbox}
\DeclareUnicodeCharacter{2192}{\ensuremath{\rightarrow}}
\DeclareUnicodeCharacter{2191}{\ensuremath{\uparrow}}
\DeclareUnicodeCharacter{2193}{\ensuremath{\downarrow}}
\DeclareUnicodeCharacter{25B2}{\ensuremath{\blacktriangle}}
\DeclareUnicodeCharacter{25BC}{\ensuremath{\blacktriangledown}}
\usepackage{amssymb}

% Palette de l'application : ivoire, bordeaux, or
\definecolor{bordeaux}{HTML}{7A1F2B}
\definecolor{bordeauxl}{HTML}{F6E7E9}
\definecolor{or}{HTML}{B8860B}
\definecolor{orl}{HTML}{FBF3DC}
\definecolor{sapin}{HTML}{2C7A55}
\definecolor{sapinl}{HTML}{E2F1E9}
\definecolor{sky}{HTML}{1C6DB4}
\definecolor{skyl}{HTML}{E6F0F8}
\definecolor{summit}{HTML}{6B7A8F}
\definecolor{night}{HTML}{2A2420}
\definecolor{ink}{HTML}{1F2430}
\AtBeginDocument{\color{ink}}

\titleformat{\section}
  {\Large\bfseries\sffamily\color{bordeaux}}
  {\makebox[1.6em][l]{\thesection}}{0.3em}{}
  [{\vspace{1pt}\color{bordeaux!35}\titlerule[1.3pt]}]
\titlespacing*{\section}{0pt}{14pt}{7pt}
\titleformat{\subsection}{\large\bfseries\sffamily\color{night}}{\thesubsection}{0.5em}{}
\pretocmd{\section}{\needspace{10\baselineskip}}{}{}
\pretocmd{\subsection}{\needspace{7\baselineskip}}{}{}
\widowpenalty=10000
\clubpenalty=10000
\raggedbottom
\setlength{\parindent}{0pt}
\setlength{\parskip}{6pt}
\setlist[enumerate]{itemsep=3pt,leftmargin=6mm}
\setlist[itemize]{itemsep=2pt,leftmargin=5mm}

\pagestyle{fancy}
\fancyhf{}
\renewcommand{\headrulewidth}{0.4pt}
\renewcommand{\headrule}{\hbox to\headwidth{\color{bordeaux!40}\leaders\hrule height \headrulewidth\hfill}}
\fancyfoot[C]{\sffamily\small\color{summit}\thepage}
\renewcommand{\sectionmark}[1]{\markboth{#1}{}}

\newcommand{\menu}[1]{\textsf{\textbf{\textcolor{bordeaux}{#1}}}}
\newcommand{\touche}[1]{\tcbox[on line,boxsep=0pt,left=2pt,right=2pt,top=1.5pt,bottom=1.5pt,
  colback=white,colframe=summit,boxrule=0.5pt,arc=1.5pt]{\footnotesize\textsf{#1}}}
\newcommand{\fichier}[1]{\texttt{\small #1}}

\newcommand{\boite}[5]{\newtcolorbox{#1}{enhanced,unbreakable,colback=#2,colframe=#3,
  coltitle=white,fonttitle=\bfseries\sffamily\small,boxrule=0.9pt,arc=2.5pt,
  left=2.5mm,right=2.5mm,top=1.2mm,bottom=1.2mm,title={#4~~#5}}}
\newcommand{\capture}[2][1]{\par\needspace{8\baselineskip}\begin{center}\setlength{\fboxsep}{0pt}\setlength{\fboxrule}{0.6pt}%
  \fcolorbox{summit!50}{white}{\includegraphics[width=\dimexpr#1\linewidth-1.3pt\relax]{#2}}\end{center}}

\usepackage[french]{babel}
\usepackage[colorlinks=true,linkcolor=bordeaux,urlcolor=bordeaux]{hyperref}
\hypersetup{pdftitle={Tutti — manuel d'utilisation},pdfauthor={},pdfcreator={Tutti}}
\boite{astucebox}{orl}{or!85!black}{\faLightbulb}{Bon à savoir}
\boite{securitebox}{sapinl}{sapin}{\faShield*}{Vos fichiers sont en sécurité}
\boite{attentionbox}{bordeauxl}{bordeaux}{\faExclamationTriangle}{Attention}
\newcommand{\astuce}[1]{\par\needspace{5\baselineskip}\begin{astucebox}#1\end{astucebox}}
\fancyhead[L]{\sffamily\small\color{bordeaux}\textbf{TUTTI — MANUEL}}
\fancyhead[R]{\sffamily\small\color{summit}\leftmark}
\input{version}

\begin{document}

\begin{titlepage}
\centering
\vspace*{2.2cm}
\includegraphics[width=3.6cm]{../../logo_1024.png}\par
\vspace{8mm}
{\sffamily\Huge\bfseries\color{bordeaux} TUTTI\par}
\vspace{3mm}
{\sffamily\LARGE\color{ink} Assembler, convertir, retoucher et découper des PDF\par}
\vspace{2mm}
{\sffamily\large\color{summit} Manuel d'utilisation \textbullet{} version \versiontutti\par}
\vspace{12mm}
\begin{tcolorbox}[enhanced,width=0.8\linewidth,colback=orl,colframe=or,arc=4pt,boxrule=1.2pt,left=5mm,right=5mm,top=3mm,bottom=3mm]
\sffamily\centering
\textit{Tutti} : « tous ensemble », en musique.\\[2mm]
\raggedright
Quatre onglets, dans l'ordre d'un travail :\\[1mm]
\textbf{Assembler} \quad → \quad \textbf{Pages} \quad → \quad \textbf{Finitions} \qquad et \quad \textbf{Découper}
\end{tcolorbox}
\vfill
{\footnotesize\color{summit}\sffamily Windows \textbullet{} macOS \textbullet{} Linux\par}
\end{titlepage}

\tableofcontents
\newpage

\section{À quoi sert Tutti}

Tutti remplace les sites Internet du genre « iLovePDF » : tout se fait sur votre ordinateur,
rien n'est envoyé sur Internet. Chaque onglet de la fenêtre fait un travail :

\begin{center}
\renewcommand{\arraystretch}{1.4}
\begin{tabularx}{\linewidth}{>{\raggedright\arraybackslash}X >{\sffamily\bfseries\color{bordeaux}}l c}
\toprule
\textbf{Je veux…} & \textbf{\color{ink}Onglet} & \textbf{Page} \\
\midrule
réunir plusieurs PDF en un seul (conducteur + parties, avec un signet par partie) & Assembler et convertir & \pageref{sec:assembler} \\
transformer une photo, un scan, un Word, un Excel ou un PowerPoint en PDF & Assembler et convertir & \pageref{sec:convertir} \\
faire un livret à plier et agrafer (A5 sur A4, ou A4 sur A3) & Assembler et convertir & \pageref{sec:options} \\
redresser une page scannée de travers ou couchée & Pages & \pageref{sec:pages} \\
retirer une page blanche ou en double, changer l'ordre des pages & Pages & \pageref{sec:pages} \\
numéroter les pages, écrire un titre en haut de chaque page & Finitions & \pageref{sec:finitions} \\
rendre un PDF assez léger pour l'envoyer par e-mail & Finitions & \pageref{sec:finitions} \\
faire un fichier par page, ou extraire des pages & Découper & \pageref{sec:decouper} \\
\bottomrule
\end{tabularx}
\end{center}

\begin{securitebox}
\begin{itemize}
\item Les fichiers d'origine ne sont \textbf{jamais modifiés} : Tutti les lit, il n'écrit que le nouveau PDF.
\item Les pages sont recopiées \textbf{telles quelles}, sans perte de qualité ni « réimpression ».
\item Tutti refuse d'écrire par-dessus l'un des fichiers à assembler.
\item Une coupure de courant pendant l'écriture ne laisse jamais de PDF à moitié écrit.
\end{itemize}
\end{securitebox}

\section{Installer Tutti}

Téléchargez la dernière version sur la page \textbf{Releases} du projet, en choisissant le fichier
de votre système :

\begin{center}
\renewcommand{\arraystretch}{1.3}
\begin{tabular}{>{\bfseries}l l}
\toprule
Windows & \fichier{Tutti-Setup-X.Y.Z.exe} \\
macOS (puces Apple M1, M2…) & \fichier{Tutti-X.Y.Z-macos-arm64.dmg} \\
macOS (processeur Intel) & \fichier{Tutti-X.Y.Z-macos-x86\_64.dmg} \\
Linux & \fichier{Tutti-X.Y.Z-linux-x86\_64.tar.gz} (ou \fichier{-arm64}) \\
\bottomrule
\end{tabular}
\end{center}

\textbf{Aucun mot de passe administrateur n'est demandé} : Tutti s'installe dans votre profil.

\subsection*{Windows}
Lancez \fichier{Tutti-Setup-X.Y.Z.exe}. Au tout premier lancement, Windows peut afficher
« Windows a protégé votre ordinateur », parce que l'installeur n'est pas signé :
cliquez sur \menu{Informations complémentaires}, puis \menu{Exécuter quand même}.
Cet avertissement n'apparaît qu'une seule fois.

\subsection*{macOS}
Ouvrez le fichier \fichier{.dmg} et faites glisser Tutti dans \menu{Applications}.
Au premier lancement, macOS refuse d'ouvrir une application non signée :
faites un \textbf{clic droit} sur Tutti, choisissez \menu{Ouvrir}, puis confirmez.

\subsection*{Linux}
Décompressez l'archive et lancez \fichier{installer.sh}. Tutti s'installe dans
\fichier{\textasciitilde/.local/share} et apparaît dans le menu des applications.

\section{Premier lancement}

Tutti s'ouvre dans la langue de votre système, français ou anglais. Le menu \menu{Langue}
permet de forcer l'une ou l'autre.

Une question vous est posée une seule fois : \textbf{envoyer automatiquement un rapport si
Tutti rencontre un problème}. Le rapport contient la version du programme et du système et
l'erreur rencontrée ; les noms de vos fichiers et votre nom d'utilisateur sont masqués.
Vous pouvez changer d'avis à tout moment dans le menu \menu{Aide}.

\capture[0.82]{../capture_vide-fr.png}

\section{Assembler des PDF, pas à pas}\label{sec:assembler}

\subsection{Ajouter les fichiers}
\begin{itemize}
\item \textbf{Glissez} les PDF depuis l'Explorateur (ou le Finder) dans la grande liste ;
\item ou glissez \textbf{un dossier entier} : tous les PDF qu'il contient sont ajoutés, rangés par nom ;
\item ou utilisez \menu{Ajouter des fichiers…} \ (\touche{Ctrl}+\touche{O}) et \menu{Ajouter un dossier…}.
\end{itemize}
Chaque fichier affiche son nombre de pages et sa taille. Un fichier déjà présent n'est pas ajouté
une seconde fois.

\subsection{Les ranger}
L'ordre de la liste est l'ordre du PDF final.
\begin{itemize}
\item \textbf{Faites glisser} un fichier vers le haut ou vers le bas dans la liste ;
\item ou sélectionnez-le et utilisez \menu{▲ Monter} / \menu{▼ Descendre}
      (\touche{Alt}+\touche{↑} / \touche{Alt}+\touche{↓}) ;
\item ou cliquez sur \menu{Trier par nom}. Le tri est « naturel » : \fichier{2 Flûte} passe avant
      \fichier{10 Violon}, comme on l'attend.
\end{itemize}
\menu{Retirer} \ (\touche{Suppr}) enlève les fichiers sélectionnés de la liste ;
\menu{Tout effacer} vide la liste. \textbf{Aucun fichier n'est supprimé du disque.}

\capture[0.82]{../capture_triee-fr.png}

\subsection{Vérifier le nom du fichier final}
Tutti propose un nom à partir de la partie commune des noms des fichiers :
\fichier{Symphonie n°5 - 01 Conducteur}, \fichier{Symphonie n°5 - 02 Flûte}… donnent
\fichier{Symphonie n°5 - complet.pdf}, dans le dossier du premier fichier. Vous pouvez le
modifier directement, ou choisir un autre emplacement avec \menu{Choisir…}.

\subsection{Assembler}
Cliquez sur \menu{Assembler}. Une barre montre l'avancement ; le bouton devient
\menu{Annuler} pendant le travail. Une fois terminé, le PDF s'ouvre (option
\menu{Ouvrir le PDF une fois terminé}).

\astuce{Si le fichier final existe déjà, Tutti demande s'il faut le remplacer. S'il est ouvert
dans un lecteur PDF, fermez-le d'abord : Windows interdit de remplacer un fichier ouvert.}

\section{Les options}\label{sec:options}

\begin{description}[style=nextline,leftmargin=5mm,font=\sffamily\bfseries\color{bordeaux}]
\item[Un signet par fichier]
Le PDF final contient un signet par fichier d'origine (« 02 Flûte », « 10 Violon I »…), affiché
dans le panneau latéral de votre lecteur PDF. Le nom du signet est celui du fichier, sans la
partie commune à tous.
\item[Chaque partie commence sur une page de droite]
Pour l'impression recto-verso : si une partie se termine sur une page de droite, une page blanche
est insérée, pour que la partie suivante commence elle aussi sur une page de droite.
\item[Mise en page : livret]
Pour un \textbf{livret plié et agrafé au milieu} (programme de concert, par exemple). Le livret
se fait à partir de \textbf{tout ce qui est dans la liste}, dans l'ordre de la liste : un seul
fichier de plusieurs pages, ou plusieurs fichiers d'une page chacun (page 1, page 2…), ou un
mélange. Le nombre de pages est complété par des pages blanches jusqu'à un multiple de 4 ; elles
se placent à la fin du livret. Trois choix :
\begin{itemize}
\item \textbf{Livret A5 prêt à imprimer} : deux pages réduites en A5, côte à côte sur chaque face
      d'une feuille A4 paysage. Imprimez en A4, paysage, \textbf{taille réelle (100 \%)},
      \textbf{recto-verso avec retournement sur le bord court}.
\item \textbf{Livret A4 sur feuilles A3} : le format des partitions reliées. Deux pages A4,
      \textbf{sans réduction}, côte à côte sur chaque face d'une feuille A3 paysage. Il faut une
      imprimante A3 (ou un imprimeur) : A3, paysage, taille réelle, recto-verso bord court.
      On obtient un cahier au format A4.
\item \textbf{Ordre du livret seulement} : les pages gardent leur format, simplement remises dans
      l'ordre du livret (8, 1, 2, 7, 6, 3, 4, 5 pour 8 pages). À l'impression, choisissez
      \textbf{2 pages par feuille}, paysage, recto-verso bord court.
\end{itemize}
Pliez ensuite la pile au milieu, emboîtez les feuilles et agrafez au pli.
\item[Ouvrir le PDF une fois terminé]
Ouvre le résultat dans votre lecteur PDF habituel. Décochée, Tutti propose d'ouvrir le PDF
ou son dossier.
\end{description}

\section{Quand un fichier pose problème}

Un fichier abîmé ou protégé par un mot de passe est signalé \textcolor{bordeaux}{en rouge} dans
la liste, avec la raison, \textbf{avant} l'assemblage. Le bouton \menu{Assembler} reste grisé
tant qu'il est présent : sélectionnez-le et cliquez sur \menu{Retirer}.

\capture[0.82]{../capture_erreur-fr.png}

\begin{center}
\renewcommand{\arraystretch}{1.35}
\begin{tabularx}{\linewidth}{>{\raggedright\arraybackslash}p{0.36\linewidth} X}
\toprule
\textbf{Message} & \textbf{Que faire} \\
\midrule
PDF illisible ou abîmé & Rouvrez le document d'origine et exportez-le à nouveau en PDF. \\
Protégé par un mot de passe & Enregistrez une copie sans mot de passe depuis votre lecteur PDF. \\
Ce PDF ne contient aucune page & Le fichier est vide : retirez-le. \\
Le fichier est sans doute ouvert dans un autre programme & Fermez le PDF final dans votre lecteur, puis recommencez. \\
\bottomrule
\end{tabularx}
\end{center}

\section{Convertir des images et des documents en PDF}\label{sec:convertir}

Dans l'onglet \menu{Assembler et convertir}, la liste accepte aussi :
\begin{center}
\renewcommand{\arraystretch}{1.3}
\begin{tabularx}{\linewidth}{>{\bfseries}l X}
\toprule
Images & JPEG, PNG, TIFF (y compris plusieurs pages), BMP, GIF, WebP \\
Word & \fichier{.doc}, \fichier{.docx}, \fichier{.odt}, \fichier{.rtf} \\
Excel & \fichier{.xls}, \fichier{.xlsx}, \fichier{.ods} \\
PowerPoint & \fichier{.ppt}, \fichier{.pptx}, \fichier{.odp} \\
\bottomrule
\end{tabularx}
\end{center}

Chaque fichier est converti dès qu'il arrive dans la liste ; il apparaît avec la mention
« (image) » ou « (document converti) » et son nombre de pages. Vous pouvez ensuite le ranger
et l'assembler avec des PDF comme n'importe quel autre fichier.

\textbf{Pour convertir un seul fichier}, glissez-le seul dans la liste : Tutti propose
le même nom avec l'extension \fichier{.pdf} (\fichier{Lettre.docx} → \fichier{Lettre.pdf}),
puis \menu{Assembler} fabrique le PDF.

\begin{itemize}
\item \textbf{Images} : une image par page, au format A4 dans le sens de l'image, sans déformation.
      Une photo JPEG est insérée telle quelle : aucune perte de qualité.
\item \textbf{Word, Excel, PowerPoint} : la mise en page est confiée à \textbf{Microsoft Office}
      s'il est installé (Windows), sinon à \textbf{LibreOffice} (gratuit, \fichier{libreoffice.org}),
      indispensable sur Mac et Linux. Le logiciel travaille invisible, en lecture seule, macros coupées.
      Si aucun des deux n'est installé, le fichier est signalé en rouge.
\end{itemize}

\section{Retoucher les pages (onglet Pages)}\label{sec:pages}

L'onglet \menu{Pages} montre \textbf{toutes les pages d'un document en vignettes}, pour les
pivoter, les retirer ou changer leur ordre. Il travaille sur \textbf{un document à la fois}.

\subsection{Ouvrir le document}
Glissez \textbf{un} fichier dans le grand cadre, ou cliquez sur \menu{Ouvrir un document…}.
Un PDF, mais aussi un Word, un PowerPoint, une image… (converti d'abord, comme dans le premier
onglet). Les vignettes apparaissent l'une après l'autre ; sous chacune, le \textbf{numéro de la
page dans le document d'origine}.

\capture[0.82]{../capture_pages-fr.png}

\subsection{Choisir des pages}
\begin{itemize}
\item un \textbf{clic} choisit une page (elle se colore) ;
\item \touche{Ctrl}+clic ajoute ou retire une page du choix ; \touche{Maj}+clic choisit toute une suite ;
\item \menu{Tout choisir} (\touche{Ctrl}+\touche{A}) choisit toutes les pages ;
\item un \textbf{double-clic} affiche la page en grand, pour la lire avant de décider.
\end{itemize}

\subsection{Pivoter, retirer, déplacer}
\begin{description}[style=nextline,leftmargin=5mm,font=\sffamily\bfseries\color{bordeaux}]
\item[↺ Pivoter à gauche, ↻ Pivoter à droite]
Un quart de tour à chaque clic (\touche{Ctrl}+\touche{L} / \touche{Ctrl}+\touche{R}), sur toutes les
pages choisies. Une page pivotée porte le signe ↻ sous sa vignette. Le contenu de la page n'est pas
touché : seule son orientation change, sans aucune perte de qualité.
\item[Retirer]
Enlève les pages choisies du résultat (\touche{Suppr}). Le document d'origine, lui, garde toutes
ses pages.
\item[Déplacer]
\textbf{Faites glisser} une vignette (ou plusieurs vignettes choisies) à sa nouvelle place.
\item[Tout annuler]
Revient aux pages d'origine, dans l'ordre d'origine.
\end{description}
La ligne sous les vignettes résume les changements : « 12 pages — 1 retirée(s) — 2 pivotée(s) —
ordre changé ».

\subsection{Enregistrer}
\menu{Enregistrer…} demande où écrire le nouveau PDF ; Tutti propose
\fichier{<nom> - retouché.pdf} à côté du document. \textbf{Le document d'origine n'est jamais
modifié} : Tutti refuse d'écrire par-dessus. Si vous fermez Tutti ou ouvrez un autre document
alors que des changements ne sont pas enregistrés, Tutti demande d'abord.

\astuce{Une partition scannée dont une page est arrivée couchée : ouvrez-la dans \menu{Pages},
cliquez la page couchée, \menu{Pivoter}, puis \menu{Enregistrer…}. Trente secondes.}

\section{Finitions : numéros, titre, poids}\label{sec:finitions}

L'onglet \menu{Finitions} ajoute la touche finale à un ou \textbf{plusieurs} PDF à la fois — par
exemple à toutes les parties d'une partition. Chaque fichier donne un \textbf{nouveau} fichier,
\fichier{<nom> - finitions.pdf}, à côté de l'original (ou dans le dossier de votre choix).

\capture[0.82]{../capture_finitions-fr.png}

\begin{description}[style=nextline,leftmargin=5mm,font=\sffamily\bfseries\color{bordeaux}]
\item[Numéros de page]
\begin{itemize}
\item \textbf{Où} : en bas au centre ; ou \textbf{côté extérieur}, c'est-à-dire à droite sur
      les pages impaires et à gauche sur les paires — le bord libre d'un recto-verso relié, là
      où le pouce tourne la page ; en bas ou en haut.
\item \textbf{Forme} : « 3 », ou « 3 / 12 » (numéro et nombre total de pages).
\item \textbf{Commencer à} : le numéro de la première page, pour une partie qui fait suite à une autre.
\item \textbf{Pas de numéro sur la 1re page} : la page de titre n'est pas numérotée, mais elle
      compte (la suivante porte le 2).
\end{itemize}
\item[Texte en haut de page]
Le titre de l'œuvre, le nom de la partie (« Smetana — La Moldau — Violon I »)… centré en haut de
\textbf{toutes les pages}, ou de la \textbf{première seulement}. Laissez la case vide pour ne rien
écrire.
\item[Alléger pour l'envoi par e-mail]
Les photos et les scans en couleur ou en gris trop fins sont réduits à 150 points par pouce —
largement assez pour l'écran et une imprimante de bureau. Le texte et les notes de musique
écrits par un logiciel (Finale, Sibelius, Word…) ne sont pas des images : ils restent
parfaitement nets. Les scans en noir et blanc sont déjà légers et restent intacts. À la fin, Tutti
affiche le poids avant et après.
\end{description}

\begin{center}
\setlength{\fboxsep}{0pt}\setlength{\fboxrule}{0.6pt}%
\fcolorbox{summit}{white}{\includegraphics[width=0.42\linewidth]{../exemple_finitions-fr.png}}\\[1mm]
{\small\sffamily\color{summit} Résultat : titre en haut, « 2 / 9 » en bas côté extérieur (page paire : à gauche).}
\end{center}

\astuce{Numéros et texte s'écrivent en petits caractères (10 points), à 7 mm du bord : au-delà de
la marge que l'imprimante ne sait pas imprimer, en général en deçà de la première portée.
Si une partition remplit sa page jusqu'au bord, préférez les numéros \og en haut\fg{} ou pas de texte.}

\section{Découper un PDF}\label{sec:decouper}

Onglet \menu{Découper} :
\begin{enumerate}
\item glissez le ou les fichiers à découper dans la liste : des PDF, mais aussi des documents
      Word, Excel, PowerPoint ou des images (convertis automatiquement, comme à l'onglet précédent) ;
\item choisissez \menu{Une page par fichier}, ou \menu{Par plages de pages} en écrivant les
      plages séparées par des virgules : \fichier{1-8, 9-16, 17-} (« \fichier{17-} » veut dire
      « de la page 17 à la fin ») ;
\item choisissez où ranger les fichiers : par défaut, un dossier \fichier{<nom> - pages} est créé
      à côté de chaque PDF ;
\item cliquez sur \menu{Découper}.
\end{enumerate}
Les fichiers s'appellent \fichier{Partition - page 01.pdf}, \fichier{Partition - page 02.pdf}… ou
\fichier{Partition - pages 01-08.pdf} : les zéros gardent le bon ordre dans l'Explorateur.
Le PDF d'origine n'est pas modifié. Si des fichiers du même nom existent déjà, Tutti demande
avant de les remplacer.

\capture[0.82]{../capture_decoupe-fr.png}

\subsection*{Combiner les fonctions}
\begin{itemize}
\item \textbf{Un Word de 6 pages en 6 PDF d'une page} : glissez le \fichier{.docx} dans l'onglet
      \menu{Découper}, choisissez \menu{Une page par fichier}.
\item \textbf{Des photos et des PDF en un seul document} : glissez-les ensemble dans l'onglet
      \menu{Assembler et convertir}, rangez-les, assemblez.
\item \textbf{Un tableau Excel en PDF} : glissez-le seul dans \menu{Assembler et convertir}.
\item \textbf{Les pages 3 à 5 d'une présentation} : glissez le \fichier{.pptx} dans
      \menu{Découper}, plages \fichier{3-5}.
\item \textbf{Une partition complète, propre et prête à envoyer} :
      \menu{Assembler} les parties → ouvrir le résultat dans \menu{Pages} pour redresser une page
      couchée et retirer une page blanche → \menu{Finitions} : numéros « 3 / 12 » côté extérieur,
      titre en haut, \og alléger\fg{} → envoyer le fichier \fichier{- finitions.pdf}.
\item \textbf{Un livret de programme de concert} : assemblez les pages (Word, images…), puis
      mise en page \menu{Livret A5} — ou, pour numéroter d'abord, \menu{Finitions} puis
      \menu{Assembler} le fichier numéroté en livret.
\end{itemize}

\section{Mises à jour}

Au démarrage, Tutti vérifie discrètement s'il existe une nouvelle version. Si c'est le cas, il
propose de l'installer : il la télécharge, l'installe et redémarre. Rien n'est à télécharger
à la main. Vous pouvez aussi vérifier vous-même : \menu{Aide} → \menu{Rechercher une mise à jour}.

\section{Signaler un problème}

\menu{Aide} → \menu{Signaler un problème…} : décrivez ce qui s'est passé, puis \menu{Envoyer}.
Sans connexion Internet, le signalement est enregistré dans un fichier que Tutti vous indique.

\section{Confidentialité}

\begin{itemize}
\item Tutti n'accède à Internet que pour vérifier les mises à jour et, avec votre accord, envoyer
      les rapports d'incident.
\item Aucun de vos fichiers n'est envoyé. Dans les rapports, les noms de fichiers et le nom
      d'utilisateur sont masqués.
\end{itemize}

\section{Raccourcis clavier}

\begin{center}
\renewcommand{\arraystretch}{1.35}
\begin{tabular}{ll}
\toprule
\multicolumn{2}{l}{\sffamily\bfseries\color{bordeaux} Assembler et convertir} \\
\touche{Ctrl}+\touche{O} & Ajouter des fichiers \\
\touche{Alt}+\touche{↑} \ / \ \touche{Alt}+\touche{↓} & Monter / descendre le fichier sélectionné \\
\touche{Suppr} & Retirer le fichier sélectionné de la liste \\
\midrule
\multicolumn{2}{l}{\sffamily\bfseries\color{bordeaux} Pages} \\
\touche{Ctrl}+\touche{L} \ / \ \touche{Ctrl}+\touche{R} & Pivoter à gauche / à droite \\
\touche{Suppr} & Retirer les pages choisies \\
\touche{Ctrl}+\touche{A} & Choisir toutes les pages \\
double-clic & Voir la page en grand \\
\midrule
\multicolumn{2}{l}{\sffamily\bfseries\color{bordeaux} Partout} \\
\touche{F1} & Ouvrir ce manuel \\
\touche{Ctrl}+\touche{Q} & Quitter \\
\bottomrule
\end{tabular}
\end{center}

\vfill
{\footnotesize\color{summit}\sffamily
Tutti utilise qpdf (via pikepdf), QtPdf pour les vignettes, et les polices Source Sans 3 et EB Garamond (SIL Open Font
License). Exemples de ce manuel : fichiers fictifs.\par}

\end{document}
